위상 정렬(Topology Sort)은 방향 비순환 그래프(Directed Acyclic Graph, DAG)의 모든 정점을 방향성에 맞게 나열하는 것을 말한다. 즉, 그래프에 존재하는 모든 간선이 앞에서 뒤로 향하도록 정점들을 나열하는 것이다. 위상 정렬은 작업의 선후 관계를 나타내는 그래프에서 작업을 수행하는 순서를 결정하는 데 유용하다. 예를 들어, 어떤 작업이 다른 작업보다 먼저 완료되어야 한다면, 그 작업을 먼저 수행하도록 작업 순서를 정하는 것이다. 이번 장에서는 방향 비순환 그래프에서 위상 정렬을 수행하는 알고리즘을 다룬다.

\section{위상 정렬}
방향 비순환 그래프는 방향 그래프 중에서 사이클이 없는 방향 그래프를 말한다. 만약 사이클이 있다면 위상 정렬은 수행될 수 없다. 사이클에 속하는 그 어떤 작업도 동일 사이클에 있는 다른 작업에 의존하기 때문에, 어떤 작업도 먼저 수행될 수 없기 때문이다. 따라서 위상 정렬은 방향 비순환 그래프에서만 정의된다.

\begin{definition}[위상 정렬]
    유향 그래프 $G=(V, E)$의 위상 정렬은 $G$의 모든 정점을 나열하는 것으로, 이때 모든 간선이 앞에서 뒤로 향하도록 정점들을 나열하는 것을 말한다.
\end{definition}

위상 정렬을 구하는 방법은 큐를 활용한 방법과 DFS를 활용한 방법이 있다. 두 알고리즘 모두 $O(V+E)$의 시간복잡도로 위상 정렬을 구할 수 있다.

\section{큐 활용 위상정렬 알고리즘}
방향 비순환 그래프의 재미있는 성질 중 하나는 진입 차수가 0인 정점이 무조건 존재한다는 것이다. 만약 그렇지 않다면 사이클이 존재하게 된다. 더 나아가, 진입 차수가 0인 정점과, 그 정점과 연결된 간선들을 제거하면 여전히 방향 비순환 그래프가 된다. 따라서 진입 차수가 0인 정점을 큐에 넣고, 큐에서 하나씩 꺼내면서  그 정점과 인접한 정점들의 간선을 제거시키는 과정을 반복하면 위상 정렬이 가능하다. 큐에서 꺼낸 순서대로 놓으면 되기 때문이다.

여기서 문제는 간선을 제거하는 것이다. 인접행렬 방식의 구현이라면 간선 제거가 매우 쉽겠지만, 인접 리스트 방식에서는 어렵다. 이러한 문제의 해결을 위해서는 그래프의 인접 리스트를 업데이트하지 않고, 단순히 진입 차수만 관리하는 배열을 만든 후, 자신으로 들어오는 간선 하나가 사라질 때마다 진입 차수를 관리하는 배열의 값을 1씩 감소시키면 된다. 그러다가 진입차수가 0이 되면 큐에 삽입하는 것이다. 구현 코드는 \Cref{code: queue-topology-sort}를 참고하라.

\begin{code} \label{code: queue-topology-sort}
큐를 활용한 위상 정렬 구현
\begin{lstlisting}
#include <bits/stdc++.h>
using namespace std;
int n, m, in[202020];
vector<int> adj[202020];
queue<int> q;
vector<int> ans;

int main() {
    // make adj
    for (int i = 1; i <= n; i++) {
        for (int v : adj[i]) in[v]++;
    }
    for (int i = 1; i <= n; i++) {
        if (in[i] == 0) q.push(i);
    }
    while (!q.empty()) {
        int v = q.front();
        q.pop();
        ans.push_back(v);
        for (int u : adj[v]) {
            in[u]--;
            if (in[u] == 0) q.push(u);
        }
    }
    // ans contains the topological order
    // if ans.size() < n, there is a cycle in the graph
}
\end{lstlisting}
\end{code}

시간복잡도는 $O(V+E)$이다.

\section{DFS 활용 위상정렬 알고리즘}
DFS를 활용한 위상 정렬 알고리즘은 DFS의 종료 시간을 활용한다. 핵심 아이디어는
\begin{obs}
    방향 비순환 그래프의 각 $v$에서 $u$로 가는 간선에 대해, $v$의 종료 시간이 $u$의 종료 시간보다 이르다.
\end{obs}
이다. 위상 정렬은 모든 간선이 앞에서 뒤로 향하도록 정점을 나열하는 것이므로, 종료 시간 순으로 정점을 나열하면 위상 정렬이 된다. 구현 코드는 \Cref{code: dfs-topology-sort}를 참고하라.

\begin{code} \label{code: dfs-topology-sort}
DFS를 활용한 위상 정렬 구현
\begin{lstlisting}
#include <bits/stdc++.h>
using namespace std;
int n, m;
vector<int> adj[202020];
int ch[202020], val[101010];
stack<int> stk;
void dfs(int v) {
    ch[v] = 1;
    for (int i = 0; i < adj[v].size(); i++) {
        int u = adj[v][i];
        if (ch[u]) continue;
        dfs(u);
    }
    stk.push(v);
}
int main() {
    // make adj
    for (int i = 1; i <= n; i++) ch[i] = 0;
    for (int i = 1; i <= n; i++) {
        if (ch[i]) continue;
        dfs(i);
    }
    while (stk.size() > 0) {
        int v = stk.top();
        stk.pop();
        // v is in topological order
    }
}
\end{lstlisting}
\end{code}

이 방법은 주어진 그래프가 DAG일 때는 위상정렬을 구해주지만, 사이클이 있는지 여부는 탐지하지 못한다는 단점이 있다.

\section{연습문제 A}
\begin{problem}
    \Cref{code: dfs-topology-sort}를 약간만 변형하여 시간복잡도를 유지하면서 사이클이 있는지 여부를 탐지할 수 있다. 주어진 그래프가 DAG인지 판별하는 코드를 추가하라.
\end{problem}
% \begin{problem}
%     최소 스패닝 트리로 가능한 것이 하나인지, 여러 개인지 판별하는 프로그램을 작성하라.
% \end{problem}
% \begin{problem}
%     $N$개의 수로 이루어진 수열 $A_i$가 주어진다. 여러분은 $i$번 정점과 $j$번 정점을 비용 $\gcd(A_{i+1}, A_{i+2}, \cdots, A_j)$로 연결할 수 있다. $1$번부터 $N$번까지의 정점이 포함된 그래프 $G$를 연결 그래프로 만들기 위해 필요한 최소 비용을 $O(N^2)$의 시간복잡도로 구하는 프로그램을 작성하라. 단, $\gcd$는 최대공약수 함수이다.
% \end{problem}

\section{연습문제 B}

\subsection{기초문제}
% \begin{problem} \url{https://www.acmicpc.net/problem/1197} \end{problem}
% \begin{problem} \url{https://www.acmicpc.net/problem/4343} \end{problem}
% \begin{problem} \url{https://www.acmicpc.net/problem/4386} \end{problem}
% \begin{problem} \url{https://www.acmicpc.net/problem/9134} \end{problem}
% \begin{problem} \url{https://www.acmicpc.net/problem/18769} \end{problem}
% \begin{problem} \url{https://www.acmicpc.net/problem/20010} \end{problem}
% \begin{problem} \url{https://www.acmicpc.net/problem/24010} \end{problem}
% \begin{problem} \url{https://www.acmicpc.net/problem/25545} \end{problem}
% \begin{problem} \url{https://www.acmicpc.net/problem/27009} \end{problem}
% \begin{problem} \url{https://www.acmicpc.net/problem/28019} \end{problem}

\subsection{응용문제}
% \begin{problem} \url{https://www.acmicpc.net/problem/1396} \end{problem}
% \begin{problem} \url{https://www.acmicpc.net/problem/3900} \end{problem}
% \begin{problem} \url{https://www.acmicpc.net/problem/7788} \end{problem}
% \begin{problem} \url{https://www.acmicpc.net/problem/10335} \end{problem}
% \begin{problem} \url{https://www.acmicpc.net/problem/12667} \end{problem}
% \begin{problem} \url{https://www.acmicpc.net/problem/14574} \end{problem}
% \begin{problem} \url{https://www.acmicpc.net/problem/18919} \end{problem}
% \begin{problem} \url{https://www.acmicpc.net/problem/20263} \end{problem}
% \begin{problem} \url{https://www.acmicpc.net/problem/23108} \end{problem}
% \begin{problem} \url{https://www.acmicpc.net/problem/23736} \end{problem}

\subsection{도전문제}
% \begin{problem} \url{https://www.acmicpc.net/problem/4203} \end{problem}
% \begin{problem} \url{https://www.acmicpc.net/problem/5257} \end{problem}
% \begin{problem} \url{https://www.acmicpc.net/problem/8527} \end{problem}
% \begin{problem} \url{https://www.acmicpc.net/problem/11592} \end{problem}
% \begin{problem} \url{https://www.acmicpc.net/problem/12668} \end{problem}
% \begin{problem} \url{https://www.acmicpc.net/problem/14828} \end{problem}
% \begin{problem} \url{https://www.acmicpc.net/problem/16074} \end{problem}
% \begin{problem} \url{https://www.acmicpc.net/problem/17582} \end{problem}
% \begin{problem} \url{https://www.acmicpc.net/problem/18482} \end{problem}
% \begin{problem} \url{https://www.acmicpc.net/problem/18963} \end{problem}
% \begin{problem} \url{https://www.acmicpc.net/problem/19251} \end{problem}
% \begin{problem} \url{https://www.acmicpc.net/problem/19852} \end{problem}
% \begin{problem} \url{https://www.acmicpc.net/problem/20086} \end{problem}
% \begin{problem} \url{https://www.acmicpc.net/problem/20503} \end{problem}
% \begin{problem} \url{https://www.acmicpc.net/problem/24953} \end{problem}
% \begin{problem} \url{https://www.acmicpc.net/problem/24961} \end{problem}
% \begin{problem} \url{https://www.acmicpc.net/problem/25121} \end{problem}
% \begin{problem} \url{https://www.acmicpc.net/problem/27071} \end{problem}
% \begin{problem} \url{https://www.acmicpc.net/problem/27877} \end{problem}
% \begin{problem} \url{https://www.acmicpc.net/problem/28002} \end{problem}